\section{A scalar normalization without squareclass corrections}
\label{sec:scalar-correction-free}

Put
\[
                         \beta(X)=216X^4+450459
\]
and retain the five-contact polynomial
\[
 r(T)=T(T-1)(T-2401)(32T+343)(2T-1).
\]
Define
\[
 r_X=r(\beta(X)),\qquad
 \delta_X=1+(17r_X)^2,\qquad
 R_X=\frac{102^2r_X}{\delta_X^2}.
\]
The scaling factor is an integer.  In particular, this normalization
introduces no fixed denominator prime.

\begin{lemma}\label{sn:integral-input}
For $X\in\mathbb Q$ one has
\[
                 \beta(X)\in\mathbb Z
                 \quad\Longleftrightarrow\quad X\in\mathbb Z.
\]
For every real $X$, $\beta(X)\ge450459>2401$, so $r_X>0$
and $\delta_X>0$.  If $p>76176$ and $v_p(\beta(X))<0$,
then
\[
 v_p(\beta(X))=4v_p(X),\qquad
 v_p(r_X)=5v_p(\beta(X)),\qquad
 v_p(R_X)=-3v_p(r_X).
\]
In particular, all three valuations are even.
\end{lemma}
\begin{proof}
If $X$ is an integer, so is $\beta(X)$.  Conversely, write
$X=a/b$ in lowest terms with $b>0$.  Integrality of $\beta(X)$
gives $b^4\mid216=2^3\cdot3^3$.  Neither prime exponent
admits a nontrivial fourth power, so $b=1$.

The real inequalities follow from the fourth power and the
five positive factors of $r(T)$ when $T>2401$.  At a prime
in the stated range, a negative valuation of $\beta$ forces
$v_p(X)<0$.  The unique dominant terms in $\beta$ and
$r(\beta)$ are respectively $216X^4$ and $64\beta^5$.
This gives the first two valuation formulas.  In this case
$v_p(r_X)<0$, so the dominant term of $\delta_X$ is
$17^2r_X^2$, and $v_p(\delta_X)=2v_p(r_X)$.  The final
formula follows.  The displayed multiple of four is even.
\end{proof}

\begin{lemma}\label{sn:local-squares}
If $X\in\mathbb Z$, then $r_X$ is a nonzero square in
$\mathbb Q_2$, $\mathbb Q_3$, and $\mathbb Q_{17}$.
More precisely,
\[
 v_2(r_X)=2,\quad r_X/4\equiv1\pmod8,
 \qquad
 v_3(r_X)=2,\quad r_X/9\equiv1\pmod3,
\]
and at seventeen either $v_{17}(r_X)=0$ with square-unit
residue, or
\[
                    v_{17}(r_X)=2,\qquad
                    r_X/17^2\equiv9\pmod{17}.
\]
\end{lemma}
\begin{proof}
Since $216$ is divisible by eight and by twenty-seven,
\[
                    \beta\equiv3\pmod8,
                    \qquad\beta\equiv18\pmod{27}.
\]
For the two even factors of $r(\beta)$ one has
\[
 v_2(\beta-1)=v_2(\beta-2401)=1,
 \qquad
 \frac{\beta-1}{2}-\frac{\beta-2401}{2}=1200\equiv0\pmod8.
\]
Their odd parts therefore have square product congruent to one
modulo eight.  The remaining odd factors have product
$3\cdot7\cdot5\equiv1\pmod8$.  This proves the assertions
at two.  At three, only the factor $\beta$ is divisible by
three; it has valuation two.  Reduction after division by nine
gives
\[
 \frac{r(\beta)}9\equiv
       2\cdot2\cdot2\cdot1\cdot2\equiv1\pmod3.
\]

At seventeen the fourth powers are exactly $0,1,4,13,16$.
Since $216\equiv12$ and $450459\equiv10\pmod{17}$, all
possibilities are given by the table
\[
\begin{array}{c|rrrrr}
 X^4&0&1&4&13&16\\\hline
 \beta&10&5&7&13&15\\
 r(\beta)&0&15&2&13&2.
\end{array}
\]
The nonzero residues are squares: $13=8^2$, $15=7^2$, and
$2=6^2$ modulo seventeen.

In the remaining case $17\mid X$.  The exact identity
\[
 32\cdot450459+343=17^2\cdot49879,
 \qquad 49879\equiv1\pmod{17},
\]
and $v_{17}(X^4)\ge4$ show that the factor
$32\beta+343=6912X^4+17^2\cdot49879$ has valuation exactly
two and unit residue one.  The other four factors are units,
with product
$10\cdot9\cdot6\cdot2\equiv9\pmod{17}$.
This proves the last assertion.  The usual local square criteria
at two and at odd primes now prove all three square statements.
\end{proof}

\begin{proposition}\label{sn:prepared-constant}
For every integer $X$, the number $R_X$ is prepared: it is
positive and less than one, its reduced denominator is a square,
it is a square in $\mathbb Q_2,\mathbb Q_3,\mathbb Q_{17}$,
and
\[
 v_2(R_X)=4,\qquad v_3(R_X)=4,\qquad
 v_{17}(R_X)\in\{2,4\}.
\]
No squareclass correction is required.
\end{proposition}
\begin{proof}
The contact value $r_X$ is a positive integer.  One has
$\gcd(r_X,\delta_X)=1$, and the preceding lemma shows that
$\delta_X$ is a unit at two and three; it is a unit at
seventeen by its definition.  Thus
$\gcd(102^2r_X,\delta_X^2)=1$, and the reduced denominator
is exactly $\delta_X^2$.

The quotient $R_X/r_X=(102/\delta_X)^2$ is a rational square.
Its local valuations add two at each of $2,3,17$, proving
both the asserted squareclasses and the exact depths.  Finally,
since $r_X\ge1$,
\[
 0<R_X=
 \frac{102^2r_X}{(1+17^2r_X^2)^2}
 <\frac{102^2}{17^4r_X^3}
 \le\frac{10404}{83521}<1.
\]
\end{proof}

\subsection{The norm tree and its domain}

Fix a positive rational number $\rho$ and six rational tests
$t_0,\ldots,t_5$, with their assigned positive integer orders.
Start the scalar norm tree with the distinguished node
$1+t_0^2/\rho$.  Its other initial nodes are the assigned
normalized tests or pair norms.  Every later operation is
\[
                  N_\rho(u,v)=u^2+\rho v^2
\]
or the unchanged carrying of an unpaired node.  This is the
constant-left tree used by the scalar certificate; denote its
root by $W_\rho$.

\begin{lemma}\label{sn:positive-root}
Where its rational expressions are defined, $W_\rho>0$.
When all six tests tend to zero, $W_\rho$ tends to one.
For fixed nonzero rational $\rho$, at any finite set of primes
the tests can be required sufficiently small that
$W_\rho$ is a unit at every prime in that set.
\end{lemma}
\begin{proof}
The distinguished initial node is positive.  At each norm
operation its positive descendant contributes a positive
square, with a positive coefficient if it is the second
argument.  Its path therefore remains positive up to the
root.  When all tests vanish, this path starts at one and
all other nodes vanish, so its root is one.  For fixed
$\rho$, the tree is a polynomial in the normalized tests
with fixed rational coefficients.  It is continuous at zero
over $\mathbb R$ and over every $\mathbb Q_p$.
Its value at zero is one.  The inverse image of $\mathbb Z_p^*$
is therefore a neighborhood of zero, proving the final claim.
\end{proof}

There are two valid implementations.  One may replace the contact
parameter throughout the tree by $\rho=R_X$ and put
$b=R_XW_{R_X}^2$.  Alternatively, one may retain the original
tree with $\rho=r_X$ and normalize only its final pole input:
\[
                       b=R_XW_{r_X}^2.
\]
The second choice retains the original test denominators.  It
also has the quadratic scaling equations
\[
        \delta_X=1+17^2r_X^2,\qquad
        \delta_X V=102W_{r_X},\qquad b=r_XV^2,
\]
where the last expression can use the existing square-and-product
gates.  The first two displayed equations are quadratic.
The change of arithmetic construction is not asserted to be an
identity on every former auxiliary solution.

For rational inputs, both implementations have $r_X,R_X,\delta_X>0$
and $W_\rho>0$.  Thus the scalar contact, normalizer, norm-root,
and pole-input nonvanishing conditions follow on the stated
rational-expression domain.  A coordinate used solely to certify
one of these nonzero conditions can be removed after its occurrence
condition is checked.  This observation does not remove unrelated
denominator guards in the elliptic-point formulas.

\subsection{Soundness and simultaneous completeness}

\begin{proposition}\label{sn:contact-preservation}
At every finite scalar contact prime $p>76176$, the replacement
has the same positive odd contact valuation as the original
certificate.  Consequently it preserves its norm-tree extraction
of the assigned contact orders, for either implementation above.
\end{proposition}
\begin{proof}
At such a contact, $\nu=v_p(r_X)>0$ is odd.  The normalizer
satisfies $\delta_X\equiv1\pmod p$, and $102$ is a unit.
Hence
\[
               v_p(R_X)=v_p(r_X)=\nu,
 \qquad R_X/r_X\in\mathbb Z_p^{*2}.
\]
For either $\rho=r_X$ or $R_X$, the two terms in
$u^2+\rho v^2$ have valuations of opposite parity when both
are nonzero.  There is therefore no cancellation: its valuation
is the minimum of $2v_p(u)$ and $\nu+2v_p(v)$.
Every step of the existing constant-left extraction uses exactly
this identity, the assigned powers of the contact parameter in
the leaves, and the odd-pole exclusion for $b$.  All these
valuations are unchanged.  Lemma~\ref{sn:integral-input}
also shows that a contact with infinity has even valuation
outside the fixed exceptional set, so it creates no omitted
odd contact.  The inherited height and comparison arguments
therefore give $\beta\in\mathbb Z$, and
Lemma~\ref{sn:integral-input} gives $X\in\mathbb Z$.
\end{proof}

\begin{proposition}\label{sn:simultaneous-completeness}
Under the same arithmetic inputs as the source construction,
the scalar block with the input map and normalization above
can be satisfied at every integer input, together with the
remaining point, height, comparison, and size conditions.
The finite squareclass correction parameter can be omitted.
\end{proposition}
\begin{proof}
Fix an integer input.  The numbers $r_X,R_X,\delta_X$ are
now fixed, before any contact primes are chosen.  Proposition~
\ref{sn:prepared-constant} supplies the sign, local squareclasses,
square denominator, and small positive real size formerly obtained
by a correction.  Include the finite support of the numerator
and denominator of $R_X$, and that of $\rho$, in the fixed
local approximation conditions for the common point $H$ and
the scalar point $S-P$.

The required approximation is the same one used in the source.
Choose $H$ sufficiently deep at these finite places and in a
small real neighborhood of the identity, excluding the identity
itself.  Shrink the real neighborhood so that its formal
parameter is bounded away from zero.  Choose $S-P$ still deeper
at every fixed finite place and sufficiently close to the real
identity.  The reciprocal-abscissa tests at these points then
tend to zero, and the quotient of their formal parameters
tends to zero with $H$ in that smaller neighborhood.
Fixed powers of $\rho$ in the denominators do not affect this
convergence.

One can impose these conditions with the required congruences
on the integer indices: the source first chooses the fixed
modulus and depths, then applies its prime-pattern and common-point
construction.  Increasing the modulus to include the finite
support of $\delta_X$ is allowed because this support was fixed
before that application.  It introduces no new family of
varying-prime conditions and no stronger prime-pattern theorem.
The real compact-group argument permits the same shrinking of
the neighborhoods and arbitrarily large scalar index.  All
other source index and comparison choices are unchanged.

By Lemma~\ref{sn:positive-root}, these choices make
$W_\rho$ a unit at every prime in the support of $R_X$ and
give $|W_\rho-1|<1/2$ at the real place.  Therefore
\[
                       0<b=R_XW_\rho^2<9/4<17.
\]
The factor $W_\rho^2$ preserves the local squareclasses and,
at $2,3,17$, preserves the depths of Proposition~
\ref{sn:prepared-constant}.  The reduced denominator of $b$
is a square: at a prime in the support of $R_X$, the root is
a unit, while at every other prime its valuation contribution
is even.  Thus $b$ is prepared.  The same inherited rational
pole-realization theorem supplies the final pole witnesses.
This proves completeness without a correction branch.
\end{proof}

The input polynomial has degree four and can be evaluated with
two quadratic gates while preserving the source expression
$\beta=2Y+3$:
\[
                    Z=X^2,\qquad Y=108Z^2+225228.
\]
The gain in any compiled polynomial also depends on its exact
occurrence and denominator inventory; the preceding propositions
establish the arithmetic substitution independently of that count.

\subsection{Number-field scope}

All displayed algebraic identities remain valid in characteristic
zero.  Positivity of the contact and norm root also holds in a
number field with a real embedding, after applying that embedding.
The local squareclass table, however, concerns rational integer
inputs and the three specified completions of $\mathbb Q$.
It does not establish the corresponding preparation assertions
for every algebraic integer in a larger field.

Even the integrality equivalence of Lemma~\ref{sn:integral-input}
does not automatically become an equivalence for rings of
integers.  Let $x_0$ be a real fourth root of $1/108$ and
$L=\mathbb Q(x_0)$.  Then
\[
                       \beta(x_0)=450461\in\mathbb Z.
\]
But $x_0$ is not an algebraic integer, since $x_0^4=1/108$.
Thus this elementary input map
alone does not transfer a definition to $\mathcal O_L$.
No assertion that the complete rational defining polynomial
defines every $\mathcal O_L$ is made or follows from this
normalization.
